\section*{Broader Impact}

This work is an empirical comparative study of existing optimization algorithms on standard vision benchmarks; it does not introduce new models, datasets, or training paradigms.
The principal practical benefit lies in computational efficiency: spectral optimizers reach near-ceiling accuracy within 20 epochs, matching or surpassing the accuracy that AdamW attains at 50 epochs.
Adopting spectral methods in place of longer AdamW schedules can reduce GPU hours by $2\times$ or more, directly lowering energy consumption and carbon emissions associated with model training.

Beyond efficiency, the fair hyperparameter tuning protocol developed here (independent per-pair learning rate sweeps with transparent reporting of sweep outcomes) provides a reusable methodology for optimizer evaluation.
Standardized comparison protocols reduce the risk of misleading benchmark results and help practitioners make informed optimizer choices without redundant experimentation.

The primary limitation of this study is dataset scale: all experiments use CIFAR-100, Tiny-ImageNet, and ImageNet-100.
Whether the observed spectral advantage and architecture-dependent patterns persist at ImageNet-1K scale or on larger models remains an open question.
Practitioners should validate these findings on their target tasks before making deployment decisions.
We identify no direct negative societal impacts from this methodological research.
